\section{Formal model and checked mathematics}
\label{sec:lean-model}

The formal development is written in Lean~4 on top of mathlib
\cite{deMouraUllrich2021Lean4,MathlibCommunity2020}.  This section gives the
mathematical core in the order used by the soundness proof: extract a spend
witness from an algebraic trace, establish the concrete evaluation codes,
carry one decoded candidate through four folds, connect those folds to the
algebraic relation, authenticate every queried value, and prove the spend
state transition.  The implementation and probability arguments consume
these theorems in \cref{sec:implementation,sec:soundness}.

\subsection{From algebraic rows to a spend witness}

The proof system opens normalized columns encoding range checks, value
balance, asset equality, ownership, note construction, nullifier construction,
public fields, and two Merkle paths.  The formal model records the residual
equation for every constraint.  Hash rows contain the intermediate states of
the Poseidon2 permutation, while each Merkle level records its sibling and the
left/right position selected by the path bit.

\begin{theorem}[Trace soundness for the spend relation]
\label{thm:spend-after-extraction}
Let \(x\) be a public statement and let \(O\) be the normalized columns opened
by the proof.  Suppose that all specified arithmetic residuals vanish, the
opened hash rows follow the specified Poseidon2 round function, both
depth-twenty paths are well formed, and the six spend fields in \(O\) equal
the corresponding fields in \(x\).  If the deployed note, nullifier, and
Merkle-node functions implement that round function on the reached inputs,
then \(x\) has a witness satisfying Definition~\ref{def:spend}.
\end{theorem}

The Lean proof constructs the witness from the opened columns.  It converts
the field balance equation into \cref{eq:value-balance} using the
\(2^{30}\) range bounds, rebuilds both twenty-level hash chains from their
intermediate rows, derives ownership and the nullifier from the input-note
opening, and obtains both note equations from the same trace.  The range
argument is essential: since each quantity is below \(2^{30}\), equality in
\(\mathbb F_{2^{31}-1}\) lifts to equality of the represented integers with
no modular wraparound.

This theorem supplies the final semantic step of extraction:
\[
  \text{valid authenticated algebraic trace}
  \Longrightarrow
  \text{valid private-spend witness}.
\]
The earlier steps establish that a successful verifier execution supplies
that trace or enters one of the failure events counted in the security
experiment.

\subsection{Field, circle domain, and encoders}

The base field is represented as \(\mathbb Z/(2^{31}-1)\mathbb Z\).  Lean
constructs the unit circle as a commutative group and proves that the chosen
generator has order \(2^{31}\).  It follows that two powers of the generator
have the same horizontal coordinate exactly when their exponents agree up to
sign.  The degree-four challenge field is built as a fixed tower over \(\F\),
including the packing basis and arithmetic identities used by the verifier.

The program stores circle evaluations in bit-reversed order and groups four
symbols into each query fibre.

\begin{theorem}[Concrete domains and encoders]
\label{thm:domains-encoders}
The selected circle generator over \(\F\) has order \(2^{31}\).  Every point
in the initial circle domain and three later line domains is distinct.  Under
bit-reversed storage, each query fibre maps to its four specified circle
positions.  All four encoders equal polynomial evaluation on their stated
domains and in their deployed order.  The later agreement bounds are
\(255,63,15,3\), which gives the required distance hypotheses.
\end{theorem}

The order proof reduces equality of horizontal coordinates to equality of
generator exponents up to sign.  The encoder proofs then unfold the exact
permutation and fibre maps used by the verifier.  This establishes more than
the cardinality of the domains: it identifies each stored symbol with the
polynomial evaluation that the fold equations require.

The published circle-code decoding result is supplied at a named mathematical
interface.  Lean verifies its application to the concrete field, domain
sizes, agreement threshold, list bound, nineteen batched columns, and maximum
batching exponent eighteen.  The general decoder theorem comes from the
literature; every parameter used to instantiate it here is checked by the
formal development.

\subsection{One candidate through four folds}

List decoding may return several nearby codewords.  Soundness requires one
initial candidate to remain coherent through all four verifier challenges;
choosing a fresh candidate after each challenge would grant the prover an
unpriced adaptive choice.  The formal model therefore indexes each prover
response by the transcript prefix available when that response is formed.

For each round \(r\in\{1,2,3,4\}\), Lean constructs a bad-challenge set
\(B_r\) fixed by the preceding prefix and proves its cardinality bound.  It
then proves the following causal extraction result.

\begin{theorem}[Coherent four-fold extraction]
\label{thm:four-fold-extraction}
An accepted ideal low-degree execution either samples a challenge from one of
the four prefix-determined sets \(B_r\), or there is one member of the initial
decoder list whose four exact radix-four folds equal the final polynomial
published in the transcript.
\end{theorem}

The extraction proof has complementary backward and forward parts.  Starting
from the four published coefficients, the backward argument selects close
predecessors at dimensions \(16\), \(64\), \(256\), and \(1024\).  Outside
the four predecessor-failure events this gives one chain
\[
  c^{(0)}\in\K^{1024},\quad
  c^{(r+1)}=F_{\alpha_r}(c^{(r)}),\quad
  c^{(4)}=c_{\mathrm{final}}.
\]
Only \(c^{(0)}\) is chosen from the bounded decoder list; the later vectors
are its folds, so the list-size factor is paid once.  For Fiat--Shamir, the
forward argument chooses a deterministic nearby candidate from each available
transcript prefix.  A transition from a state with no coherent continuation
to one with a continuation can occur only when the next challenge lies in a
bad set already fixed by that prefix.  This is the causal statement needed by
the oracle reduction.

The same four challenges update the weights used by the algebraic relation.
Let \(F_\alpha:\K^{4n}\to\K^n\) be the natural coefficient fold, let
\(D_\alpha:\K^{4n}\to\K^n\) be the verifier's dual weight fold, and let
\(\partial\) denote the boundary functional on the seven coefficients of a
round polynomial.

\begin{proposition}[Candidate and weight fold duality]
\label{prop:fold-duality}
For candidate values \(v\in\K^{4n}\) and verifier weights
\(w\in\K^{4n}\), the deployed convolution constructs a polynomial
\(P_{v,w}\) of degree at most six satisfying
\[
  \partial P_{v,w}=\sum_{i=0}^{4n-1}v_iw_i,
  \qquad
  P_{v,w}(\alpha)
  =\sum_{j=0}^{n-1}F_\alpha(v)_jD_\alpha(w)_j.
\]
These identities hold at dimensions \(1024\), \(256\), \(64\), and \(16\)
with the coefficient and weight orders used by the Rust verifier.
\end{proposition}

For one four-element fibre, the coefficient fold uses weights
\((1,\alpha,\alpha^2,\alpha^3)\), while the dual fold uses
\((1,\alpha^3,\alpha^2,\alpha)/4\).  Expanding their product gives the seven
coefficients of \(P_{v,w}\).  Summing this identity over fibres proves the
proposition.  The relation argument applies it in every round and bounds the
challenge tuples capable of repairing a false scalar claim.  Prover messages
may depend on earlier challenges, while each repair set is fixed before the
challenge charged to it.

\subsection{Byte-level transcript}

The transcript is an ordered sequence of typed operations.  An absorb event
contains a label and exact byte payload; a squeeze event records the requested
field element, point, selector, or query block; a work event records its nonce
and difficulty.  These operations expand to primitive SHA-256 calls.  An
absorb advances the transcript state once, a squeeze uses separate output and
state-advance hashes, and a work check examines its nonce without advancing
the state.

Lean proves the complete operation order, event counts, labels, payload
lengths, unique positions, and the check-before-absorb order for all six work
nonces.  It also proves that the challenges later consumed by the low-degree
and relation checks are precisely the values produced at their transcript
positions.  The translated Rust path decodes the fixed proof-body slices into
the same transcript inputs; \cref{sec:implementation} connects that decoder to
the formal schedule.

\subsection{Authenticated openings}

Five Merkle commitments authenticate the queried values.  An opened record is
the value bytes followed by a 32-byte public salt.  Leaf, binary-node, and
radix-four-node hashes have separate domains, include the tree identity, and
preserve child order.  The two initial trees use query index \(q\); the three
later trees use \(q\!\gg\!2\), \(q\!\gg\!4\), and \(q\!\gg\!6\).
Sorted, duplicate-free indices determine the compressed frontier.  Each
odd-depth path ends in a binary cap selected by the remaining index bit.

For each tree, the exact parser trace records every opened record, path slot,
frontier node, and unconsumed suffix.  The five-tree routine feeds each suffix
into the following parser and requires the final suffix to be empty.  From
these traces, Lean constructs an authenticated forest and proves that every
value consumed by the low-degree checker is the value portion of one of its
authenticated records.

\begin{theorem}[Authenticated consumption]
\label{thm:authenticated-consumption}
Fix the five roots and the eighteen query positions.  If the opening routine
succeeds, every low-degree value read by that execution belongs to the exact
record authenticated for its tree and position.  Two successful forests that
expose different reached values under those roots yield a concrete SHA-256
collision witness.
\end{theorem}

This collision-witness formulation avoids an injectivity assumption for a
compressing hash function.

\subsection{State semantics}

The state model assigns five accounts to fixed roles and checks their keys,
owners, writable flags, current pool root and sequence, marker availability,
and derived marker address.  On success it writes the nullifier marker,
updates the root, increments the sequence, and returns the exact resulting
state.  A pre-existing marker causes rejection, so two sequential modeled
successes cannot use the same nullifier.  Rejection leaves the modeled visible
state unchanged.

A separate cleanup transition closes the proof account, transfers its full
modeled balance to the supplied refund account, and preserves the pool and
nullifier state.  The deployed interpretation of account locking, rollback,
and finalized persistence is isolated in the runtime assumptions of
\cref{sec:limitations}.

\begin{table}[t]
\centering
\small
\caption{Main machine-checked mathematical results and their role.}
\label{tab:formal-results}
\begin{tabularx}{\textwidth}{@{}R{0.22\textwidth}YY@{}}
\toprule
Result & Checked content & Use in the security argument \\
\midrule
Spend extraction & Arithmetic, ownership, notes, nullifier, public fields,
and both Merkle paths imply a witness for Definition~\ref{def:spend}. & Turns a valid
extracted trace into a valid private spend. \\
Domains and encoders & Circle-generator order, distinct evaluation points,
stored ordering, encoder identities, and later-code distances. & Discharges
the concrete hypotheses of the decoding argument. \\
Four-fold extraction & One coherent initial candidate reaches the final
polynomial, or a prefix-timed bad challenge occurs. & Prevents adaptive
switching between decoder candidates. \\
Fold duality & The degree-six round polynomial carries the current candidate
dot product to the candidate and weight folds. & Joins low degree to the
algebraic relation. \\
Transcript correctness & Exact byte payloads, operation order, work placement,
and challenge consumption. & Supplies the causal random-coin schedule. \\
Opening correctness & Exact five-tree parsing and authentication of every
value consumed by the low-degree checker. & Binds checked values to committed
proof data. \\
State correctness & Exact successful marker and pool update, replay rejection,
rollback in the state model, and cleanup preservation. & Gives a one-time
atomic ledger transition. \\
\bottomrule
\end{tabularx}
\end{table}
